\section{Problem Formulation}
\label{sec:problem}

\subsection{Visibility and Dwell}

At time step $k$, target $i$ is visible from chaser position $r_k$ when range,
field-of-view, incidence-angle, and line-of-sight predicates are satisfied:
\begin{align}
  \rho_{\min} \le \norm{p_i-r_k} \le \rho_{\max}, \label{eq:range}\\
  p_i \in \fov(r_k,q_k), \label{eq:fov}\\
  \arccos\left(
    \frac{(r_k-p_i)^{\mathsf{T}}n_i}
    {\norm{r_k-p_i}\norm{n_i}}
  \right) \le \theta_{\max}, \label{eq:view-angle}\\
  \los(r_k,p_i,\mathcal{O}) = 1, \label{eq:los}
\end{align}
where $q_k$ denotes camera attitude or an equivalent boresight model,
$\theta_{\max}$ is the maximum view angle, and $\mathcal{O}$ is the occluding
geometry representation. The dwell update is
\begin{equation}
  \tau_{i,k+1} =
  \begin{cases}
    \tau_{i,k}+\dt, & \text{if target } i \text{ is visible at } k,\\
    \tau_{i,k}, & \text{otherwise}.
  \end{cases}
\end{equation}

\subsection{Safety and Actuation}

The chaser must satisfy acceleration, speed, and clearance constraints:
\begin{align}
  \norm{u_k}_2 &\le u_{\max}, \\
  \norm{v_k}_2 &\le v_{\max}, \\
  d_{\mathcal{S}}(r_k) &\ge d_{\mathrm{safe}},
\end{align}
where $d_{\mathcal{S}}(r_k)$ is the signed distance from the chaser to the
station safety envelope. In the implemented ROS stack, a safety monitor logs
minimum distance, clearance, caution-zone status, nearest primitive, filter
activity, and nominal versus filtered acceleration commands.

\subsection{Optimization Objective}

The ideal finite-horizon inspection problem is
\begin{align}
  \min_{\{x_k,u_k,q_k\}_{k=0}^{K}}
  \quad &
  w_t K\dt
  + w_u \sum_{k=0}^{K-1} \norm{u_k}_2\dt
  + w_s \sum_{k=0}^{K} \phi(d_{\mathcal{S}}(r_k))
  + w_c (1-C_K) \label{eq:objective}\\
  \mathrm{s.t.}\quad &
  x_{k+1}=f_{\mathrm{CW}}(x_k,u_k,\dt), \\
  & \norm{u_k}_2 \le u_{\max}, \\
  & \norm{v_k}_2 \le v_{\max}, \\
  & d_{\mathcal{S}}(r_k) \ge d_{\mathrm{safe}}, \\
  & y_{i,k} \text{ follows the visibility and dwell model},\\
  & C_K \ge C_{\mathrm{req}},
\end{align}
where $\phi$ is a soft penalty used in relaxed formulations. This problem is
mixed-integer and nonconvex because visibility, occlusion, dwell satisfaction,
and coverage are discrete or discontinuous. OrbInspect therefore implements a
structured decomposition rather than a monolithic nonlinear mixed-integer
program.

\subsection{Evaluation Metrics}

The primary metrics are final coverage, cumulative control effort, minimum
clearance, maximum speed, and tracking error:
\begin{align}
  C_{\mathrm{final}} &= \frac{N_{\mathrm{inspected}}}{N},\\
  \Delta v_{\mathrm{cum}} &= \sum_k \norm{u_k}_2 \dt,\\
  d_{\min} &= \min_k d_{\mathcal{S}}(r_k),\\
  e_{\mathrm{track,rms}} &=
    \sqrt{\frac{1}{K}\sum_k \norm{r_k-r_k^{\mathrm{ref}}}_2^2}.
\end{align}
Additional reported metrics include planning time, selected-viewpoint count,
candidate-viewpoint count, trajectory sample count, and safety-filter
activation.

\subsection{Dynamics-Aware Coverage Efficiency}

The role of the CW model is not limited to post hoc feasibility checking.  In
the implemented offline planner, each candidate viewpoint is evaluated through
a predicted bounded-acceleration CW transfer from the current state.  This
produces an edge cost
\begin{equation}
  J_{\mathrm{CW}}(x,z_j) =
  \lambda_v \widehat{\Delta v}(x,z_j)
  + \lambda_e \widehat{e}_{\mathrm{track}}(x,z_j)
  + \lambda_s \max\{0,-\widehat{d}_{\min}(x,z_j)\},
\end{equation}
where $\widehat{\Delta v}$ is the integrated acceleration norm,
$\widehat{e}_{\mathrm{track}}$ is the RMS transfer tracking error, and
$\widehat{d}_{\min}$ is the predicted minimum clearance.  The planner therefore
prefers viewpoints that expose new surface area while requiring low control
effort under orbital relative dynamics.  The reported
$\Delta v/C_{\mathrm{final}}$ metric is used as a direct measure of coverage
efficiency.
